\documentclass[12pt]{article}
\usepackage{amsmath,amssymb}
\begin{document}
\section*{Research notes}

Problem: strip \(\Sigma=\mathbb R\times[0,1]\), deck map
\(G_\beta(x,y)=(x+\beta,1-y)\).  A realized \(\beta\) gives an embedding
of the quotient \(M_\beta=\Sigma/\langle G_\beta\rangle\), a flat
\(1\times\beta\) Mobius band.  The invariant locally finite clean
triangulation descends to a finite clean triangulation.

Important definitions used in the proof:
\begin{itemize}
\item In a clean triangle, the boundary edge is the ridge and the opposite
vertex is the apex.
\item A bend is a segment from the apex to a point of the ridge.
\end{itemize}

Key local lemma: if an affine map on a clean triangle shortens the ridge
and lengthens the two side edges, then it lengthens every bend.  Proof:
for ridge endpoints \(A,B\), apex \(C\), \(P_s=(1-s)A+sB\),
\[
q(s)=|L(P_s)-L(C)|^2-|P_s-C|^2
\]
is concave because
\(q''=2(|L(A)-L(B)|^2-|A-B|^2)<0\).  Since \(q(0),q(1)>0\),
concavity gives \(q(s)>0\).

The bends form a circle naturally parametrized by the centerline
\(\{y=1/2\}/G_\beta\).  Each bend meets this centerline once and every
centerline point lies in a unique bend.  Going once around applies
\(G_\beta\), so a continuously chosen bottom-to-top orientation reverses.

T-pair lemma: for the image segments of bends, choose unit direction
\(e(u)\) and midpoint \(m(u)\), with \(e(u+1)=-e(u)\), \(m(u+1)=m(u)\).
For \(0<r<1\),
\[
\Phi(u,r)=\big(e(u)\cdot e(u+r),
(m(u)-m(u+r))\cdot(e(u)\times e(u+r))\big).
\]
The cylinder compactifies to \(S^2\) with boundary values \((1,0)\) and
\((-1,0)\).  The swap involution is antipodal and makes \(\Phi\) odd, so
Borsuk--Ulam gives a zero.  At a zero the supporting lines are
perpendicular and coplanar, hence intersect.

Lower-bound geometry: choose the two bends, label them \(T,B\) so the
intersection \(O\) of supporting lines is not in the relative interior of
\(B^*\).  Lift \(T\) from \((0,0)\) to \((t,1)\), \(0\le t<\beta\); if the displacement had absolute value at least \(\beta\), the lift would meet its \(G_\beta\)-translate.  Cut
along \(T\); boundary arcs have lengths \(H=\beta-t\), \(D=\beta+t\).
Boundary shortening and bend lengthening give
\(H>\sqrt{1+t^2}\).  An endpoint of \(B^*\) is at distance \(>1\) from
the line of \(T^*\).  The boundary arc \(A\in\{H,D\}\) containing this
endpoint has length \(A>\sqrt{5+t^2}\).  If \(A=H\), get
\(2\beta>2\sqrt{5+t^2}>2\sqrt3\).  If \(A=D\), get
\[
2\beta>\sqrt{1+t^2}+\sqrt{5+t^2},\quad
2\beta>2\sqrt{5+t^2}-2t.
\]
The first function is increasing, the second decreasing, and both equal
\(2\sqrt3\) at \(t=1/\sqrt3\), so \(\beta>\sqrt3\).

Upper bound strategy: use the classical triangular construction that for
each \(\lambda>\sqrt3\) there is an embedded polygonal isometric paper
Mobius band \(P:M_\lambda\to\mathbb R^3\) affine on a clean
triangulation.  For a target \(\beta>\sqrt3\), pick
\(\lambda\in(\sqrt3,\beta)\), set \(s=\lambda/\beta\), pull the
triangulation back by \(h(x,y)=(sx,y)\), and scale the range by
\(r\in(1,1/s)\) close to \(1/s\).  Boundary edges are multiplied by
\(rs<1\); a diagonal with \(M_\beta\)-offset \(d\) has image length
\(r\sqrt{1+s^2d^2}\), which is \(>\sqrt{1+d^2}\) for all finitely many
diagonal edge orbits if \(r\) is close enough to \(1/s\).

Computation check: for
\(a(t)=\sqrt{1+t^2}+\sqrt{5+t^2}\) and
\(b(t)=2\sqrt{5+t^2}-2t\), \(a'(t)>0\), \(b'(t)<0\), and at
\(t_0=1/\sqrt3\) both equal \(2\sqrt3\).
\end{document}
